% TOP_PROBLEM: {"id": 30003388, "title": "Sharp Convergence Domains for Cluster Expansions", "category_id": 16, "category": {"id": 16, "name": "physics", "display_name": "Mathematical Physics", "description": "Problems at the intersection of mathematics and physics.", "slug": "physics", "order_index": 16, "created_at": "2026-07-31T15:26:25.671Z"}, "status": "open", "statusQualification": "The graph and partition-function notation were recovered from PDF pages 15--17. The report already proves a sufficient Dobrushin condition: some $r:V\\to[0,1)$ satisfies $R(v)\\leq r(v)\\prod_{v'\\in N(v)}(1-r(v'))$. The inherited question about exact and sharp domains is broader than this established sufficient bound; it is not an assertion that the report leaves all convergence unexplained."}
% FIRST_POSTED: 2026-10-11
% ENTRY_KIND: statement_only
% SUBMISSION: {"kind": "statement_only", "category": "top_problems", "problem_id": "30003388", "proposal_type": "established", "source": "Mini-Workshop: Cluster Expansions: From Combinatorics to Analysis through Probability (2017), Oberwolfach Reports, DOI 10.4171/owr/2017/8, pp. 417--452. Primary report accessed 2026-10-11; relevant PDF pages 15, 16, 17. https://doi.org/10.4171/owr/2017/8; Oberwolfach report (catalogue-listed context; not independently consulted in this preparation). https://ems.press/content/serial-article-files/46675; Bissacot--Fernández--Procacci (catalogue-listed context; not independently consulted in this preparation). https://doi.org/10.1007/s10955-010-9956-1; Fernández--Nguyen (catalogue-listed context; not independently consulted in this preparation). https://arxiv.org/abs/2112.13134", "status_evidence": "Canonical source status open retained. Supplied literature triage dated 2026-08-21: The extracted question concerns the exact convergence diameter of a specific cluster-expansion series. Later work gives stronger sufficient convergence criteria, but no exact diameter for the stated series is established in the checked literature. Primary-report access checked 2026-10-11; this historical report and inherited triage do not certify current resolution of every component.", "status_checked_on": "2026-10-11", "duplicate_check": "Compared canonical ID, title and normalized target with the complete local UnsolvedMath registry; normalized exact-text matches: none; title matches: none. Related or overlapping canonical records: none identified in this batch. Primary source context and the existing OWR statement audit used for scope and duplicate review; an exhaustive mathematical-equivalence/literature search was not performed.", "attribution": "Original proposer not named in the supplied canonical record; attribution remains to the identified Oberwolfach contribution/report, not its editors as invented proposers. Explanatory statement preparation and source/semantic audit assisted by OpenAI GPT-6 Codex; no human mathematical authorship or independent proof verification is claimed."}
% SOURCE_STATEMENT: "For a polymer-model cluster expansion\n$$\n\\log Z_L(w)=\\sum_X a(X)w^X,\n$$\ndetermine its exact domain of convergence and sharp conditions on the contour weights and graph-dependent radii that guarantee convergence."
% Preparation and final self-audit: CONTRIBUTING.md prompts C and D; source/audit artifacts retained in the local physics_statements_2026_10_11 batch.
% Imported from: physics_statements.tex; UnsolvedMath ID 30003388
% !TeX program = lualatex
\documentclass[11pt]{article}
\usepackage[margin=25mm]{geometry}
\usepackage{fontspec}
\setmainfont{Latin Modern Roman}
\usepackage{amsmath,amssymb,mathtools}
\usepackage{unicode-math}
\setmathfont{Latin Modern Math}
\usepackage{xurl,hyperref}
\hypersetup{colorlinks=true,urlcolor=blue}
\setcounter{secnumdepth}{0}
\setlength{\parindent}{0pt}
\setlength{\parskip}{5pt}
\setlength{\emergencystretch}{3em}
\newcommand{\sref}[2]{\hyperlink{source:#1:#2}{[S#2]}}
\newcommand{\cataloguescope}[1]{\paragraph{Recorded target.}#1}
\newcommand{\R}{\mathbb R}
\newcommand{\C}{\mathbb C}
\newcommand{\N}{\mathbb N}
\newcommand{\Z}{\mathbb Z}
\newcommand{\Q}{\mathbb Q}
\newcommand{\F}{\mathbb F}
\newcommand{\eps}{\varepsilon}
\begin{document}
% UnsolvedMath ID 30003388; source catalogue code OWR-15213-001
\setcounter{section}{30003387}
\section{Sharp Convergence Domains for Cluster Expansions}\label{30003388}
{\small UnsolvedMath ID 30003388 · Mathematical Physics}
\subsection{Definitions and mathematical statement}

\paragraph{Shared notation.}$\mathbb N=\{1,2,\ldots\}$; the standard number fields have their usual meanings. All additional conventions are stated locally.


\paragraph{Definitions and conventions.}
A polymer model has objects called polymers, a compatibility relation and complex weights $w$. Its partition function sums products of weights over compatible configurations. A cluster expansion expresses $\log Z_L$ as a formal sum over connected incompatibility clusters; $X$ indexes a multiset and $w^X$ its weight monomial. A convergence domain is a region of the weight space in which this series converges in the specified sense. Contour weights and graph radii determine the estimates for a particular model. More explicitly, the report has an incompatibility graph $G$ on polymers $V$, finite $L\subset V$, and $Z_L(w)=\sum_{I\subset L\text{ independent}}\prod_{v\in I}w(v)$. Its graph-dependent polydisc is $\{w:|w(v)|\leq R(v),\ v\in L\}$.

\paragraph{Target retained from the canonical source catalogue.}
\paragraph{Catalogue wording retained for traceability.}
For a polymer-model cluster expansion
$$
\log Z_L(w)=\sum_X a(X)w^X,
$$
determine its exact domain of convergence and sharp conditions on the contour weights and graph-dependent radii that guarantee convergence.

\paragraph{Source scope and qualification.}
The graph and partition-function notation were recovered from PDF pages 15--17. The report already proves a sufficient Dobrushin condition: some $r:V\to[0,1)$ satisfies $R(v)\leq r(v)\prod_{v'\in N(v)}(1-r(v'))$. The inherited question about exact and sharp domains is broader than this established sufficient bound; it is not an assertion that the report leaves all convergence unexplained.
\subsection{Short English statement}
Find the exact convergence region and sharp usable convergence criteria for polymer cluster expansions.
\subsection{Sources}
\begin{itemize}
\item[\hypertarget{source:30003388:1}{S1}] ULAM.AI and the UnsolvedMath Contributors, UnsolvedMath, record 30003388 (OWR-15213-001). Catalogue curation under CC BY 4.0; definitions and English exposition added locally.
\url{https://huggingface.co/datasets/ulamai/UnsolvedMath}.
\item[\hypertarget{source:30003388:2}{S2}] Mini-Workshop: Cluster Expansions: From Combinatorics to Analysis through Probability (2017), Oberwolfach Reports, DOI 10.4171/owr/2017/8, pp. 417--452. Primary report accessed 2026-10-11; relevant PDF pages 15, 16, 17.
\url{https://doi.org/10.4171/owr/2017/8}.
\item[\hypertarget{source:30003388:3}{S3}] Oberwolfach report (catalogue-listed context; not independently consulted in this preparation).
\url{https://ems.press/content/serial-article-files/46675}.
\item[\hypertarget{source:30003388:4}{S4}] Bissacot--Fernández--Procacci (catalogue-listed context; not independently consulted in this preparation).
\url{https://doi.org/10.1007/s10955-010-9956-1}.
\item[\hypertarget{source:30003388:5}{S5}] Fernández--Nguyen (catalogue-listed context; not independently consulted in this preparation).
\url{https://arxiv.org/abs/2112.13134}.
\end{itemize}
\label{end:30003388}
\end{document}
